% Part: computability
% Chapter: recursive-functions
% Section: general-recursive-functions

\documentclass[../../../include/open-logic-section]{subfiles}

\begin{document}

\olfileid{cmp}{rec}{gen}
\olsection{الدوال العودية العامة}

ثمة طريقة أخرى للحصول على مجموعة من الدوال الكلية. لنسم الدالة الكلية
$f(x,\vec z)$ \emph{منتظمة} إذا وجد، لكل متتالية من الأعداد الطبيعية
$\vec z$، عدد $x$ بحيث $f(x,\vec z)=0$. وبعبارة أخرى، الدوال المنتظمة
هي بالضبط الدوال التي يمكن تطبيق البحث غير المحدود عليها والخروج بدالة
كلية. ويمكن، على نحو محافظ، قصر البحث غير المحدود على الدوال المنتظمة:

\begin{defn}
\ollabel{defn:general-recursive}
مجموعة \emph{الدوال العودية العامة} هي أصغر مجموعة من الدوال من
الأعداد الطبيعية إلى الأعداد الطبيعية، بمختلف الرتب، تحتوي الصفر
واللاحق والإسقاطات، وتكون منغلقة تحت التركيب والعودية البدائية والبحث
غير المحدود المطبق على الدوال \emph{المنتظمة}.
\end{defn}

من الواضح أن كل دالة عودية عامة كلية. والفرق بين
\olref{defn:general-recursive} و\olref[par]{defn:recursive-fn} هو أنه
يسمح في الثانية باستعمال دوال عودية جزئية في الخطوات الوسيطة، والشرط
الوحيد هو أن تكون الدالة التي تنتهي إليها كلية. ولذلك فكلمة «عامة»،
وهي بقايا تسمية تاريخية، في غير موضعها؛ إذ يبدو
\olref{defn:general-recursive} في ظاهره \emph{أقل} عمومًا من
\olref[par]{defn:recursive-fn}. لكن لحسن الحظ أن الفرق ظاهري؛ فمع
اختلاف التعريفين، مجموعة الدوال العودية العامة ومجموعة الدوال العودية
مجموعة واحدة بعينها.

\end{document}
